% !TEX root = ../main.tex

\section{Analytical Solution for a Monopole} \label{sec:monopole_derivation}

Analytical approaches for monopole sources have been validated in previous work by Char and Thomas \cite{TingChar2024, Thomas2012}. This section describes the derivation of the analytical monopole to provide a background on the analytical solution provided in the validation cases. Monopoles or thickness sources are sounds that are equivalent to mass supply, described as the following,

\begin{equation}
    \Box^2p'(\bm{x},t) = \frac{\partial}{\partial t} (Q(t) \delta(\bm{x} - y(t)))
    \label{eq:monopole_def}
\end{equation}

Convolving the above with the Green's function in 3D free-space (see Appendix \ref{app:GF}, described as,

\begin{equation}
    G=\frac{\delta(g(\tau))}{4\pi r}
    \label{eq:GreensFunction}
\end{equation}

where $g(\tau)=t-\tau-\frac{r(\tau)}{c}$ and $r(\tau)=|\bm{x}-y(\tau)|$, gives the following,

\begin{equation}
    p'(\bm{x}, t)= \frac{\partial}{\partial t} \left[ \frac{Q(\tau)}{4\pi r (1-M_r)}\right]_{\tau*}
    \label{eq:monopole_convolved}
\end{equation}

which is the description of the monopole used in this work.


\subsection{Monopole Pressure Equation} \label{subsec:mono_p}
Taking the monopole definition as described in Equation \ref{eq:monopole_convolved}, taking the differentiation using the definitions as described in Appendix~\ref{app:derivatives}, the analytical solution for pressure is obtained,% and using the quotient rule to differentiate each term, is obtained,

% \begin{equation*}
%     p'(\bm{x}, t) = \frac{1}{4\pi} \left[ 
%     \frac{\partial Q(\tau) / \partial \tau}{r(1-M_r)^2} -
%     \frac{Q(\tau) v_r}{r^2 |1-M_r|^2} +
%     \left(
%     - \left(\frac{\partial M}{\partial \tau} \right) \frac{Q(\tau)}{r|1-M_r|^3} +
%     \frac{Q(\tau)c(|M|^2 - M_r^2}{r^2 |1-M_r|^3}
%     \right)
%     \right]_{\tau *}
%     \label{eq:monopole_p_der_step1}
% \end{equation*}

% The second term was algebraically manipulated by multiplying both the numerator and denominator by $1 - M_r$ to establish a common denominator and substitute the velocity in the radiation direction, $v_r$, with the speed of sound term, $c$. Then, using further algebraic manipulation for terms with $r^2(1-M_r)^3$ denominators, the analytical solution for pressure is obtained,

\begin{equation}
    p'(\bm{x}, t) = \frac{1}{4\pi} \left[ 
    \frac{\partial Q(\tau) / \partial \tau}{r(1-M_r)^2} +
    \left(\frac{\partial M}{\partial \tau} \right)_r \frac{Q(\tau)}{r(1-M_r)^3} +
    \frac{Q(\tau) c (M_r - |M|^2)}{r^2(1-M_r)^3}
    \right]_{\tau*}
    \label{eq:monopole_p_der}
\end{equation}

The first term on the right-hand side describes the far-field propagation of the monopole due to its strength's temporal gradient, $\partial Q/\partial\tau$, the second term describes the far-field propagation due to acceleration, $\partial M/\partial \tau$, and the third term describes the near-field propagation.

\subsection{Monopole Velocity Equation} \label{subsec:monopole_u}

Using the momentum equation,
\begin{equation}
    \rho_0 \frac{\partial u_i}{\partial t} +
    \rho_0 \left( u_i \frac{\partial u_i}{\partial x_j} \right) =
    -\frac{\partial p}{\partial x_i} +
    \frac{\partial}{\partial x_j} \left(
    \mu \left(\frac{\partial u_i}{\partial x_j} + \frac{\partial u_j}{\partial x_i} \right)
    \right)
    \label{eq:momentum_eq}
\end{equation}

Linearising the equation and assuming no mean flow, substituting the monopole definition into the pressure term, and convolving with Green's function to replace $\partial/\partial x_i$ terms with $\partial/\partial t$ (see Appendix \ref{app:ddx_replace}), the solution for velocity is obtained, %the second terms of both left- and right-hand are zero. Defining the pressure gradient by defining the pressure as monopole pressure definition,

% \begin{equation*}
%     - \frac{\partial p}{\partial x_i} =
%     \frac{\partial}{\partial x_i} \frac{\partial}{\partial t} \left[
%     \frac{Q(\tau)}{4\pi r (1-M_r)}
%     \right]_{\tau*}
% \end{equation*}

% and cancelling the unsteady term from both sides gives the following for the momentum equation,

% % \begin{equation*}
% %     \rho_0 \frac{\partial u_i}{\partial t} =
% %     \frac{\partial}{\partial x_i} \frac{\partial}{\partial t} \left[
% %     \frac{Q(\tau)}{4\pi r (1-M_r)}
% %     \right]_{\tau*}
% % \end{equation*}

% % cancelling the unsteady term from both sides leaves,

% \begin{equation}
%     u_i = \frac{1}{\rho_0} \frac{\partial}{\partial x_i} \left[ 
%     \frac{Q(\tau)}{4\pi r (1-M_r)}
%     \right]_{\tau *}
%     \label{eq:monopole_spatial_derivative}
% \end{equation}

% Convolving with Green's function in free 3D-space is done to replace the spatial derivative into a temporal derivative,

% \begin{equation*}
%     p'(\bm{x}, t) = \frac{\partial}{\partial x_i} (Q_i \star G) =
%     - \frac{\partial}{\partial t} \left( 
%     Q_i \star \frac{x_iG}{c|\bm{x}|}
%     \right) -
%     Q_i \star \frac{x_i G}{|\bm{x}|^2}
%     \label{eq:convolve_monopole}
% \end{equation*}

% which has the same form as the integral shown in Appendix \ref{app:derivatives}, Equation \ref{eq:C5}. Substituting the monopole source terms into $Q_i$, and converting to retarded time using the equation found in Appendix \ref{app:ddx_replace} Equation \ref{eq:retarded_time_def}, gives,

% \begin{equation*}
%     \frac{\partial}{\partial x_i} \left[ 
%     \frac{Q(\tau)}{4\pi r (1-M_r)}
%     \right]_{\tau*}
%     = \frac{1}{1-M_r} \frac{\partial}{\partial\tau} \left[ 
%     \frac{Q(\tau)r_i}{4\pi c \bm{r}^2 (1-M_r)^2}
%     \right]_{\tau *}
%     + \left[ 
%     \frac{Q(\tau)r_i}{4 \pi r^3 (1-M_r)}
%     \right]_{\tau*}
%     \label{eq:monopole_velocity_underived}
% \end{equation*}

% which can then be solved analytically to give the solution for velocity,

\begin{equation}
    \begin{split}
    u_i' 
    &= \frac{1}{4\pi\rho_0} \Bigg[ \frac{(\partial Q(\tau)/\partial \tau)r_i}{c r^2(1 - M_r)^2} 
    + \frac{Q(\tau)r_i}{r^3(1-M_r)^3}
    + \left(\frac{dM}{d\tau} \right) \frac{Q(\tau)r_i}{cr^2(1-M_r)^3} \\
    &\phantom{= \frac{1}{4\pi\rho_0}
    \Bigg[} - \frac{Q(\tau)M^2r_i}{r^3(1-M_r)^3}
    - \frac{Q(\tau)M_i}{r^2(1-M_r)^2}
    \Bigg]_{\tau*}
    \end{split}
    \label{eq:monopole_velocity}
\end{equation}

where the first and third terms correspond to the far-field velocity perturbation components due to the strength and Mach number temporal gradients, respectively; and the second, fourth, and fifth terms correspond to the near-field perturbation components.

\subsection{Monopole Density Equation} \label{subsec:mono_rho}
The expression for density is based on the assumption of a linear, isentropic flow. Therefore, the expression for density is,

\begin{equation}
    \rho'=\frac{p'}{c^2}
    \label{eq:density}
\end{equation}

which is valid for both monopole and dipole sources. 